\documentclass[11pt]{article}
\usepackage[margin=1in]{geometry}
\usepackage{amsmath,amssymb,booktabs}
\setlength{\parindent}{0pt}
\begin{document}
\section*{HW08 Question 13}

\textbf{Problem.} For stretching and injury, the observed table is:

\begin{center}
\begin{tabular}{lrr}
\toprule
Result & Stretched & Not stretched\\
\midrule
Injury & 23 & 25\\
No injury & 222 & 197\\
\bottomrule
\end{tabular}
\end{center}

\textbf{Solution.}
Row totals are \(48\) and \(419\). Column totals are \(245\) and \(222\). Grand total is \(467\).

Expected counts use
\[
E_{ij}=\frac{(\text{row total})(\text{column total})}{467}.
\]
\[
E_{\text{injury, stretched}}=\frac{48(245)}{467}=25.18,\quad
E_{\text{injury, not}}=\frac{48(222)}{467}=22.82.
\]
\[
E_{\text{no injury, stretched}}=\frac{419(245)}{467}=219.82,\quad
E_{\text{no injury, not}}=\frac{419(222)}{467}=199.18.
\]

\[
\boxed{
\begin{array}{c|cc}
 & \text{Stretched} & \text{Not stretched}\\ \hline
\text{Injury} & 25.18 & 22.82\\
\text{No injury} & 219.82 & 199.18
\end{array}}
\]
\end{document}
